\begin{center}
	\zihao{4}\heiti
	%
	\LARGE{ 中国科学院大学\\ $\boldsymbol{2024}$年招收攻读硕士学位研究生入学统一考试试卷\\
	科目名称：电动力学（回忆版）}

\end{center}
\noindent
考生须知：\\
1.本试卷满分为150分,全部考试时间总计180分钟。\\
2.所有答案必须写在答题纸上,写在试题纸上或草稿纸上一律无效。

\vspace*{2ex}
\hrule
\vspace*{4ex}

\fancypagestyle{mypagestyle}{%
	\fancyhf{}% Clear header/footer
	\fancyfoot[OR]{第~\thepage~页，共2页}%
	\fancyfoot[EL]{第~\thepage~页，共2页}%
	\fancyfoot[OL]{科目名称：电动力学 }%
	\fancyfoot[ER]{科目名称：电动力学 }%
	\renewcommand{\headrulewidth}{0pt}% Header rule of .4pt
	\renewcommand{\footrulewidth}{0.5pt}% Header rule of .4pt
}
\fancypagestyle{plain}{%
	\fancyhf{}% Clear header/footer
	\fancyfoot[OR]{第~\thepage~页，共2页}%
	\fancyfoot[EL]{第~\thepage~页，共2页}%
	\fancyfoot[OL]{科目名称：电动力学 }%
	\fancyfoot[ER]{科目名称：电动力学 }%
}

\setcounter{page}{1}%设置页面计数器
\pagestyle{mypagestyle}

\onehalfspacing
% 排版结束，下面是题目
\vspace{0.4cm}

\noindent
一、选择题（共7题，每小题5分）
\vspace{0.2cm}

\noindent
1.\hspace{0.8em}以下矢量运算正确的是
\vspace{0.2cm}

A.\hspace{0.8em}$\nabla \cdot (\nabla \varphi) = \nabla ^2 \varphi $ \hspace{2em}
B.\hspace{0.8em}$\nabla( \nabla \cdot \Vec{A}) = \nabla ^2 \Vec{A} $\hspace{2em}
C.\hspace{0.8em}$\nabla \times (\nabla \varphi) = 0 $\hspace{2em}
D.\hspace{0.8em}$\nabla \cdot (\nabla \times \Vec{A}) = 0 $
\vspace{0.2cm}

\noindent
2.\hspace{0.8em}有一半径为$R$的导体球，其电荷面密度为$\sigma$，则总能量为
\vspace{0.3cm}

A.\hspace{0.8em} $\displaystyle{ \frac{4 \pi R^3 \sigma^2}{ \epsilon_0 } }$\hspace{2em} B.\hspace{0.8em} $\displaystyle{ \frac{2 \pi R^3 \sigma^2}{ \epsilon_0 } }$\hspace{2em}
C.\hspace{0.8em} $\displaystyle{ \frac{ \pi R^3 \sigma^2}{ 2 \epsilon_0 } }$\hspace{2em} 
D.\hspace{0.8em} $\displaystyle{ \frac{4 \pi R^3 \sigma}{ \epsilon_0 } }$\hspace{2em} 

\vspace{0.3cm}
\noindent
3.\hspace{0.8em}关于平面电磁波，下列说法错误的是？
\vspace{0.2cm}

A.\hspace{0.8em} $\Vec{E} \cdot \Vec{B} = 0, \quad \Vec{E} \cdot \Vec{k} = 0 , \quad \Vec{B} \cdot \Vec{k} = 0$
\vspace{0.2cm}

B.\hspace{0.8em} $\Vec{E} \cdot \Vec{k} = \Vec{B}$
\vspace{0.2cm}

 
C.\hspace{0.8em} 电场与磁场能流密度相等
\vspace{0.2cm}

D.\hspace{0.8em} 能流密度为$\Vec{E} \times \Vec{B}$
\vspace{0.4cm}

\noindent
4.\hspace{0.8em}中微子在静止时寿命为$\tau$，当它以$v/2$运动时，寿命为
\vspace{0.3cm}

A.\hspace{0.8em}$\displaystyle{ \frac{\tau}{2} }$ \hspace{2em}  
B.\hspace{0.8em}$\displaystyle{ \frac{2}{\sqrt{2}} \tau }$ \hspace{2em}  C.\hspace{0.8em}$\sqrt{3} \tau $ \hspace{2em}  
D.\hspace{0.8em}$2 \tau$
\vspace{0.3cm}

\noindent
5.\hspace{0.8em}在均匀介质内部，体极化电荷密度$\rho_p$与体自由电荷密度$\rho_f$的关系为
\vspace{0.3cm}


A.\hspace{0.8em}$\rho_p = - (1 - \displaystyle{ \frac{\varepsilon_0}{\varepsilon} } )\rho_f $\hspace{5em}
B.\hspace{0.8em}$\rho_p =  (1 - \displaystyle{ \frac{\varepsilon}{\varepsilon_0} } )\rho_f $\vspace{0.4cm}

C.\hspace{0.8em}$\rho_f = - (1 - \displaystyle{ \frac{\varepsilon_0}{\varepsilon} } )\rho_p $\hspace{5em}
D.\hspace{0.8em}$\rho_f =  (1 - \displaystyle{ \frac{\varepsilon}{\varepsilon_0} } )\rho_p $\vspace{0.3cm}

\noindent
6.\hspace{0.8em}对于超导体，以下说法错误的是？\vspace{0.2cm}

A.\hspace{0.8em}电阻为0 \hspace{2em} B.\hspace{0.8em} 内部$B$ = 0 \hspace{2em} C.\hspace{0.8em}内部$H$ = 0  \hspace{2em} D.\hspace{0.8em} 内部有超导电流$J_s$
\vspace{0.2cm}

\noindent
7.\hspace{0.8em}立方体谐振腔截止频率最低的波模为？
\vspace{0.2cm}

A.\hspace{0.8em}（0，0，0）\hspace{2em} B.\hspace{0.8em}（1，0，0）\hspace{2em} C.\hspace{0.8em}（1，1，0）\hspace{2em} D.\hspace{0.8em}（1，1，1）\vspace{0.2cm}

\noindent
8.\hspace{0.8em}
关于辐射说法正确的是？
\vspace{0.2cm}

A.\hspace{0.8em}
无论电偶极是磁偶极，场正比于
\vspace{0.2cm}

B.\hspace{0.8em}
任何运动的带电粒子都能产生辐射
\vspace{0.2cm}

C.\hspace{0.8em}
一定时，正比于，P为辐射功率
\vspace{0.2cm}

D.\hspace{0.8em}
也许是正确的一个答案

\newpage

\noindent
二、简答题（共20分）
\vspace{0.2cm}

\noindent
1.\hspace{0.8em}请简要解释梯度，旋度，散度的物理意义
\vspace{0.5cm}

\noindent
2.\hspace{0.8em}什么是极化？体极化电荷密度是什么？面极化电荷密度是什么？写出$\rho_p$与$\Vec{P}$的关系

\vspace{0.5cm}

\noindent
3.\hspace{0.8em}什么是镜像法？镜像法使用的两个条件是什么？
\vspace{0.5cm}

\noindent
4.\hspace{0.8em}请简述平面简谐波的特点
\vspace{1cm}



\noindent
三、\hspace{0.2em}某一粒子的动能W为145 MeV，$m_0c^2$为17 MeV，求$v/c$
\vspace{3cm}

\noindent
四、\hspace{0.2em}在平行板电容器内有两层介质，它们的厚度分别为$l_1$和$l_2$，电容率为$\varepsilon_1$和$\varepsilon_2$，今在两板接上电动势为$E$的电池，求；(课后题1.11)
\vspace{0.2cm}

\noindent
1、电容器两极板上的自由电荷面密度
\vspace{0.1cm}

\noindent
2、介质分界面上的自由电荷面密度\vspace{0.1cm}

\noindent
若介质是漏电的，电导率分别为$\sigma_1$和$\sigma_2$，当电流达到恒定时，上述两问题的结果如何？
\vspace{4cm}


\noindent
五、\hspace{0.2em}一对无限大的平行理想导体板，相距为$b$，电磁波沿平行于板面的$z$方向传播，设波在$x$方向上是均匀的，求可能传播的波模和每种波模的截止频率(课后习题4.14)
\vspace{4cm}

\noindent
六、\hspace{0.2em}某一带电粒子e做半径为a的非相对论性圆周运动，回旋频率为$\omega$，求(1)位置与时间的关系 (2)求出电偶极距 (3)远处的辐射电磁场，辐射能流和辐射功率。(与课后题5.11相似)